%%%%%
%
\section{Concentration massique � partir de tableau de concentrations en nombre}

\itemstyle{Op\'erateur~:} \textsf{number\_to\_mass\_concentration}
\index{number\_to\_mass\_concentration}

%%%%%OBJET%%%%
\itemstyle{Objet~:}

Calcul d'une concentration massique de particules (exprim�s en $�$g/m$^{3}$),
� partir d'un tableau de concentrations partielles en nombre (nb/cm$^{3}$)
et d'un tableau des diam�tres medians (typiquement, pour les
PCASP/SPP200, FSSP, OPC, ..).

%%%%%ARGUMENTS%%%%
\bigskip
%%%%%ENTREES%%%%
\begin{tabarg}
\multicolumn{3}{l}{\titretab {Entr\'ees}}\\
\hline
$C$ & Vecteur, \texttt{float[dimension quelconque]} & Tableau des concentrations exprim�es en nombre (nombre/cm$^{3}$)\\
$D_c$ & Coeff., \texttt{Int} & nombre de classes de $C$ \\
$\phi$ & Coeff., \texttt{float[m�me dimension que $C$]} & Tableau des diam�tres medians (micron)\\
$\rho_p$ & Coeff., \texttt{float} & Masse volumique des particules (g/cm$^3$)\\
\end{tabarg}

\medskip
%%%%%SORTIES%%%%
\begin{tabarg}
\multicolumn{3}{l}{\titretab {Sorties}}\\
\hline
$Cm$ & Vecteur, \texttt{float} & Concentration massique des particules ($�$g/m$^3$)\\
\end{tabarg}


%%%%%FORMULE%%%%
\itemstyle{Formule}

%% number_to_mass_concentration
\[
Cm = \rho_p . \frac{\pi}{6} .
\displaystyle\sum_{k=1}^{D_c} \phi_k^3.C_k
\]



%%%%%REFERENCE%%%%
\itemstyle{Commentaire~:}

Ce calcul assimile volume moyen et volume calcul� � partir du diam�tre m�dian.
Il est g�n�ralement admis que l'impr�cision induite est faible devant
les incertitudes g�n�rales de ce type de mesure.

%%%%%AUTEUR%%%%
\itemstyle{Auteur~:}
CNRM/GMEI/TRAMM

%%%%%REFERENCE%%%%
\itemstyle{R\'ef\'erences~:}
